\documentclass[12pt]{article}
\usepackage[margin=1in]{geometry}
\usepackage{setspace}
\onehalfspacing{}

% Start of preamble
%==========================================================================================%
% Required to support mathematical unicode
\usepackage[warnunknown, fasterrors, mathletters]{ucs}
\usepackage[utf8x]{inputenc}
%\usepackage{R:/sty/quiver}

\usepackage[dvipsnames,table,xcdraw]{xcolor} % colors
\usepackage{hyperref} % links
\hypersetup{
colorlinks=true,
linkcolor=blue,
filecolor=magenta,
urlcolor=cyan,
pdfpagemode=FullScreen
}

% Standard mathematical typesetting packages
\usepackage{amsmath,amssymb,amscd,amsthm,amsxtra, pxfonts}
\usepackage{mathtools,mathrsfs,xparse}

% Symbol and utility packages
\usepackage{cancel, textcomp}
\usepackage[mathscr]{euscript}
\usepackage[nointegrals]{wasysym}
\usepackage{apacite}

% Extras
\usepackage{physics}  % Lots of useful shortcuts and macros
\usepackage{tikz-cd}  % For drawing commutative diagrams easily
\usepackage{microtype}  % Minature font tweaks
%\usepackage{pgfplots} % plots

\usepackage{enumitem}
\usepackage{titling}

\usepackage{graphicx}

%\usepackage{quiver}

% Fancy theorems due to @intuitively on discord
\usepackage{mdframed}
\newmdtheoremenv[
backgroundcolor=NavyBlue!30,
linewidth=2pt,
linecolor=NavyBlue,
topline=false,
bottomline=false,
rightline=false,
innertopmargin=10pt,
innerbottommargin=10pt,
innerrightmargin=10pt,
innerleftmargin=10pt,
skipabove=\baselineskip,
skipbelow=\baselineskip]{mytheorem}{Theorem}

\newenvironment{theorem}{\begin{mytheorem}}{\end{mytheorem}}

\newtheorem{corollary}{Corollary}
\newtheorem{lemma}{Lemma}

\newtheoremstyle{definitionstyle}
{\topsep}%
{\topsep}%
{}%
{}%
{\bfseries}%
{.}%
{.5em}%
{}%
\theoremstyle{definitionstyle}
\newmdtheoremenv[
backgroundcolor=Violet!30,
linewidth=2pt,
linecolor=Violet,
topline=false,
bottomline=false,
rightline=false,
innertopmargin=10pt,
innerbottommargin=10pt,
innerrightmargin=10pt,
innerleftmargin=10pt,
skipabove=\baselineskip,
skipbelow=\baselineskip,
]{mydef}{Definition}
\newenvironment{definition}{\begin{mydef}}{\end{mydef}}

\newtheorem*{remark}{Remark}

\newtheorem*{example}{Example}
\newtheorem*{claim}{Claim}

% Common shortcuts
\def\mbb#1{\mathbb{#1}}
\def\mfk#1{\mathfrak{#1}}

\def\bN{\mbb{N}}
\def\C{\mbb{C}}
\def\R{\mbb{R}}
\def\bQ{\mbb{Q}}
\def\bZ{\mbb{Z}}
\def\cph{\varphi}
\renewcommand{\th}{\theta}
\def\ve{\varepsilon}
\newcommand{\mg}[1]{\| #1 \|}

% Often helpful macros
\newcommand{\floor}[1]{\left\lfloor#1\right\rfloor}
\newcommand{\ceil}[1]{\left\lceil#1\right\rceil}
\renewcommand{\qed}{\hfill\qedsymbol}
\renewcommand{\ip}[1]{\langle#1\rangle}
\newcommand{\seq}[2]{\qty(#1_#2)_{#2=1}^{\infty}}

\newcommand{\SET}[1]{\Set{\mskip-\medmuskip #1 \mskip-\medmuskip}}

% End of preamble
%==========================================================================================%

% Start of commands specific to this file
%==========================================================================================%

\usepackage{braket}
\newcommand{\Z}{\mbb Z}
\newcommand{\gen}[1]{\left\langle #1 \right\rangle}
\newcommand{\nsg}{\trianglelefteq}
\newcommand{\F}{\mbb F}
\newcommand{\Aut}{\mathrm{Aut}}
\newcommand{\sepdeg}[1]{[#1]_{\mathrm{sep}}}
\newcommand{\Q}{\mbb Q}
\newcommand{\Gal}{\mathrm{Gal}\qty}
\newcommand{\id}{\mathrm{id}}
\newcommand{\Hom}{\mathrm{Hom}}
\newcommand{\Ext}{\mathrm{Ext}}

%==========================================================================================%
% End of commands specific to this file

\title{Math Template}
\date{\today}
\author{Rohan Mukherjee}

\begin{document}
    \maketitle
    \begin{theorem}
        Let $0 \to \mathscr{A} \overset{\alpha}{\to} \mathscr{B} \overset{\beta}{\to} \mathscr{C} \to 0$ be a short exact sequence of cochain complexes, where for simplicity the cochain maps for $\mathscr{A}, \mathscr{B}, \mathscr{C}$ are all denoted by the same $d$. 
        \begin{enumerate}[label=(\alph*)]
            \item If $c \in C^n$ represents the class $x \in H^n(\mathscr{C})$, show that there is some $b \in B^n$ such that $\beta_n(b) = c$.
            \item Show that $d_{n+1}(b) \in \ker \beta_{n+1}$ and conclude that there is a unique $a \in A^{n+1}$ such that $\alpha_{n+1}(a) = d_{n+1}(b)$.
            \item Show that $d_{n+2}(a) = 0$ and conclude that $a$ defines a class $\overline a$ in the quotient group $H^{n+1}(\mathscr A)$.
            \item Prove that $\overline a$ is independent of the choice of $b$, i.e. if $b'$ is another choice and $a'$ is its unique preimage in $A^{n+1}$, then $\overline a = \overline{a'}$, and that $\overline a$ is also independent of the choice of $c$ representing the class $x$.
            \item Define $\delta_n(x) = \overline a$ and prove that $\delta_n$ is a group homomorphism from $H^n(\mathscr C)$ to $H^{n+1}(\mathscr A)$.
        \end{enumerate}
    \end{theorem}

    \begin{enumerate}[label=(\alph*)]
        \item By definition of an exact sequence of cochain complexes, we know that the row
        \[\begin{tikzcd}
            0 & {A^n} & {B^n} & {C^n} & 0
            \arrow[from=1-1, to=1-2]
            \arrow["{\alpha_n}", from=1-2, to=1-3]
            \arrow["{\beta_n}", from=1-3, to=1-4]
            \arrow[from=1-4, to=1-5]
        \end{tikzcd}\]
        is exact. Thus $\beta_n$ is onto, so there is some $b \in B^n$ so that $\beta_n(b) = c$.
        \item The following diagram is commutative:
        \[\begin{tikzcd}
            & \vdots & \vdots & \vdots \\
            0 & {A^n} & {B^n} & {C^n} & 0 \\
            0 & {A^{n+1}} & {B^{n+1}} & {C^{n+1}} & 0 \\
            & \vdots & \vdots & \vdots
            \arrow[from=1-2, to=2-2]
            \arrow[from=1-3, to=2-3]
            \arrow[from=1-4, to=2-4]
            \arrow[from=2-1, to=2-2]
            \arrow["{\alpha_n}", from=2-2, to=2-3]
            \arrow["{d_{n+1}}"', from=2-2, to=3-2]
            \arrow["{\beta_n}", from=2-3, to=2-4]
            \arrow["{d_{n+1}}"', from=2-3, to=3-3]
            \arrow[from=2-4, to=2-5]
            \arrow["{d_{n+1}}"', from=2-4, to=3-4]
            \arrow[from=3-1, to=3-2]
            \arrow["{\alpha_{n+1}}", from=3-2, to=3-3]
            \arrow[from=3-2, to=4-2]
            \arrow["{\beta_{n+1}}", from=3-3, to=3-4]
            \arrow[from=3-3, to=4-3]
            \arrow[from=3-4, to=3-5]
            \arrow[from=3-4, to=4-4]
        \end{tikzcd}\]
        Recall that $H^n(\mathscr C) \coloneqq \frac{\ker d_{n+1}}{\Im d_n}$. In particular, $c \in \ker d_{n+1}$ from the last step. It follows by commutativity after a slight abuse of notation that $d_{n+1}(\beta_n(b)) = d_{n+1}(c) = 0 = \beta_{n+1}(d_{n+1}(b))$, so $d_{n+1}(b) \in \ker \beta_{n+1}$. Since $\Im \alpha_{n+1} = \ker \beta_{n+1}$, there is an $a \in A^{n+1}$ so that $\alpha_{n+1}(a) = d_{n+1}(b)$. Since $\alpha_{n+1}$ is injective this $a$ is unique.

        \item We draw the third diagram:
        \[\begin{tikzcd}
            0 & {A^n} & {B^n} & {C^n} & 0 \\
            0 & {A^{n+1}} & {B^{n+1}} & {C^{n+1}} & 0 \\
            0 & {A^{n+2}} & {B^{n+2}} & {C^{n+2}} & 0
            \arrow[from=1-1, to=1-2]
            \arrow[from=1-2, to=1-3]
            \arrow[from=1-2, to=2-2]
            \arrow[from=1-3, to=1-4]
            \arrow["{d_{n+1}}"', from=1-3, to=2-3]
            \arrow[from=1-4, to=1-5]
            \arrow[from=1-4, to=2-4]
            \arrow[from=2-1, to=2-2]
            \arrow["{\alpha_{n+1}}", from=2-2, to=2-3]
            \arrow["{d_{n+2}}"', from=2-2, to=3-2]
            \arrow[from=2-3, to=2-4]
            \arrow["{d_{n+2}}"', from=2-3, to=3-3]
            \arrow[from=2-4, to=2-5]
            \arrow[from=2-4, to=3-4]
            \arrow[from=3-1, to=3-2]
            \arrow["{\alpha_{n+2}}", from=3-2, to=3-3]
            \arrow[from=3-3, to=3-4]
            \arrow[from=3-4, to=3-5]
        \end{tikzcd}\]
        We know that $\alpha_{n+1}(a) = d_{n+1}(b)$. Thus $\alpha_{n+2}(d_{n+2}(a)) = (d_{n+2} \circ d_{n+1})(b) = 0$ since $\mathscr{B}$ is a chain complex. As $\alpha_{n+2}$ is injective, $d_{n+2}(a) = 0$. Thus $\overline a$ defines a class in $H^{n+1}(\mathscr A) = \frac{\ker d_{n+2}}{\Im d_{n+1}}$, since $a \in \ker d_{n+2}$.
        \item Let $b' \in B^n$ be another element of $B^n$ with $\beta_n(b') = c$. Then $\beta_n(b - b') = 0$, so $b - b' \in \Im \alpha_n$. So find $a \in A^n$ so that $\alpha_n(a) = b - b'$. Writing $a$ and $a'$ as the preimages of $d_{n+1}(b)$ and $d_{n+1}(b')$ respectively, we can see that $\alpha_{n+1}(a - a') = d_{n+1}(b - b') = d_{n+1}(\alpha_n(a)) = \alpha_{n+1}(d_{n+1}(a))$. Since $\alpha_{n+1}$ is injective, we have that $a - a' = d_{n+1}(a)$, so $\overline{a} = \overline{a'}$.
        
        Let $y \in \Im d_n$. Since the following diagram is commutative:
        \[\begin{tikzcd}
            0 & {A^{n-1}} & {B^{n-1}} & {C^{n-1}} & 0 \\
            0 & {A^n} & {B^n} & {C^n} & 0 \\
            0 & {A^{n+1}} & {B^{n+1}} & {C^{n+1}} & 0
            \arrow[from=1-1, to=1-2]
            \arrow[from=1-2, to=1-3]
            \arrow[from=1-2, to=2-2]
            \arrow["{\beta_{n-1}}", from=1-3, to=1-4]
            \arrow["{d_n}", from=1-3, to=2-3]
            \arrow[from=1-4, to=1-5]
            \arrow["{d_n}", from=1-4, to=2-4]
            \arrow[from=2-1, to=2-2]
            \arrow[from=2-2, to=2-3]
            \arrow[from=2-2, to=3-2]
            \arrow["{\beta_n}", from=2-3, to=2-4]
            \arrow["{d_{n+1}}"', from=2-3, to=3-3]
            \arrow[from=2-4, to=2-5]
            \arrow[from=2-4, to=3-4]
            \arrow[from=3-1, to=3-2]
            \arrow["{\alpha_{n+1}}", from=3-2, to=3-3]
            \arrow[from=3-3, to=3-4]
            \arrow[from=3-4, to=3-5]
        \end{tikzcd}\]
        We know that there is some $b$ so that $\beta_{n-1}(b) = y'$. Then $\beta_n \circ d_n(b) = y$, so we can take $d_n(b)$ as the preimage of $y$ in part (b) since the class doesn't depend on the choice of $b$. Since $d_{n+1} \circ d_n(b) = 0$ and $\alpha_{n+1}$ is injective, part (b) will yield $a = 0$ and hence the class $\overline a$ corresponding to $y$ is just 0. 

        Now, for $c, c' \in C^n$, find $b$ and $b'$ so that $\beta_n(b) = c$ and $\beta_n(b') = c'$. We lift these each to $a, a'$ satisfying $\alpha_{n+1}(a) = d_{n+1}(b)$ and $\alpha_{n+1}(a') = d_{n+1}(b')$. We see then that $\alpha_{n+1}(a+a') = d_{n+1}(b+b')$, and that $\beta_n(b+b') = c+c'$, so $\overline{a + a'}$ is the class corresponding to $c + c'$. In the case where $c' \in \Im d_n$, we showed in the last paragraph that $\overline{a'} = 0$, which shows that $\overline{a}$ is independent of the choice of $c$ representing the class $x$ since every representative of $x$ can be written as the sum of $c$ and something in $\Im d_n$.

        \item We need only show that $\delta_n(x+y) = \delta_n(x) + \delta_n(y)$. Taking representatives $c, c'$ of $x,y$ respectively, if we write the class corresponding to $c, c'$ as $\overline a$ and $\overline{a'}$ resp., the last paragraph shows that the class corresponding to $c+c'$ is precisely $\overline a + \overline{a'}$. Since the $\delta_n(x)$ is independent of the choice of $c$, this shows that $\delta_n(x+y)=\delta_n(x) + \delta_n(y)$.
    \end{enumerate}


    \begin{theorem}
        Suppose $G$ is an infinite cyclic group with generator $\sigma$. 
        \begin{enumerate}[label=(\alph*)]
            \item Prove that multiplication by $\sigma-1 \in \Z G$ defines a free $G$-module resolution of 
            \[\begin{tikzcd}
                {\Z: 0} & {\Z G} & {\Z G} & {\Z } & 0
                \arrow[from=1-1, to=1-2]
                \arrow["{\sigma-1}", from=1-2, to=1-3]
                \arrow[from=1-3, to=1-4]
                \arrow[from=1-4, to=1-5]
            \end{tikzcd}\]
            \item Show that $H^0(G,A) \cong A^G$, that $H^1(G, A) \cong A/(\sigma-1)A$, and that $H^n(G, A) = 0$ for $n \geq 2$. Deduce that $H^1(G, \Z G) \cong \Z$.
        \end{enumerate}
    \end{theorem}
    
    \begin{enumerate}[label=(\alph*)]
        \item Since $\Z G$ is a free abelian group, we define a homomorphism $\pi: \Z G \to \Z$ by sending $\sigma \to 1$. We claim this yields the following projective resolution:
        \[\begin{tikzcd}
            0 & {\Z G} & {\Z G} & \Z & 0
            \arrow[from=1-1, to=1-2]
            \arrow["{\sigma-1}", from=1-2, to=1-3]
            \arrow["\pi", from=1-3, to=1-4]
            \arrow[from=1-4, to=1-5]
        \end{tikzcd}\]
        We claim that $p(x) \in \Z[x]$ has a factor of $(x-1)$ iff $p(1) = 0$. The forward direction is clear, and by applying the division algorithm to the Euclidean Domain $\Q[x]$, and using Gauss's lemma, we get a reduction of $p(x) = (x-1)q(x)$ for some $q(x) \in \Z[x]$. Thus, if $\sum_{i=0}^n a_i \sigma^i \mapsto 0$, we know that $\sum a_i = 0$, so thinking about $p(\sigma) = \sum a_i\sigma^i$ as a polynomial we write $p(\sigma) = (\sigma-1)q(\sigma)$ for some $q \in \Z[x]$. Thus $\ker \pi = (\sigma-1)\Z G$. Now, notice that if 
        \begin{align*}
            (\sigma-1)\sum_{i=0}^n a_i\sigma^i = (\sigma-1)\sum_{i=0}^m b_i\sigma^i \\
        \end{align*}
        This would translate to a linear relation on the linearly independent elements $\sigma^i$, a contradiction (where we crucially use that $|\sigma|$ is infinite). Thus we deduce $n=m$ and $a_i = b_i$ for all $0 \leq i \leq n$. So the sequence is exact, yielding a projection $G$-module resolution of $\Z$.

        \item Recall that $H^n(G, A) = \Ext^n_{\Z G}(\Z, A)$ and that $\Ext^0_{\Z G}(\Z, A) = \Hom_{\Z G}(\Z, A)$. Recall that every abelian group homomorphism $\cph: \Z \to A$ is determined by $\cph(1) = a$. Then $\cph$ is a $G$-module homomorphism iff $a = \cph(1) = \cph(\sigma^i \cdot 1) = \sigma^i \cdot a$ iff $a$ is fixed by every element of $G$. Thus $H^0(G, A) \cong A^G$ via the isomorphism $\cph \mapsto \cph(1)$. We have the following cochain complex:
        \[\begin{tikzcd}
            0 & {\Hom_{\Z G}(\Z, A)} & {\Hom_{\Z G}(\Z G, A)} & {\Hom_{\Z G}(\Z G, A)} & 0
            \arrow[from=1-1, to=1-2]
            \arrow[from=1-2, to=1-3]
            \arrow[from=1-3, to=1-4]
            \arrow[from=1-4, to=1-5]
        \end{tikzcd}\]
        We see that $\ve$ is the map sending $\cph \in \Hom_{\Z G}(\Z, A)$ to $\cph \circ \pi: \Z G \to A$, and similarly $d_1$ is the map sending $\psi \in \Hom_{\Z G}(\Z G, A)$ to $\psi \circ (\sigma-1)$, $d_2$ is the map sending $\Hom_{\Z G}(\Z G, A)$ to $0$, and $d_n$ is the identity map $0 \to 0$ for $n \geq 3$.

        We prove the following lemma:
        \begin{lemma}
            Let $\cph: A \to B$ be a homomorphism of $R$-modules. Then for any submodule $M \subset A$, there is an induced homomorphism $\overline \cph$ making the following diagram commute:
            \[\begin{tikzcd}
                A & B \\
                {A/M} & {B/\cph(M)}
                \arrow["\cph", from=1-1, to=1-2]
                \arrow[from=1-1, to=2-1]
                \arrow[from=1-2, to=2-2]
                \arrow["{\widetilde \cph}", from=2-1, to=2-2]
            \end{tikzcd}\]
            if $\cph$ is surjective, so is $\widetilde \cph$, and similarly if $\cph$ is injective, so is $\widetilde \cph$.
        \end{lemma}
        \begin{proof}
            We seek to define the map $\widetilde \cph(\overline a) = \overline{\cph(a)}$. This is well-defined since if $\overline a = \overline{a'}$, then $a - a' \in M$, so $\cph(a) - \cph(a') \in \cph(M)$, so $\overline{\cph(a)} = \overline{\cph(a')}$. $\widetilde \cph$ is obviously a homomorphism, so suppose that $\cph$ is surjective. Then for any $\overline b \in B/\cph(M)$, we can find $a \in A$ so that $\cph(a) = b$, then $\widetilde \cph(\overline a) = \overline{\cph(a)} = \overline b$, so $\widetilde \cph$ is surjective. If $\cph$ is injective, if $\widetilde \cph(\overline a) = \widetilde \cph(\overline{a'})$, then $\overline{\cph(a)} = \overline{\cph(a')}$, so $\cph(a) - \cph(a') \in \cph(M)$, so $a - a' \in M$ by injectivity, and $\overline a = \overline{a'}$, so $\widetilde \cph$ is injective.
        \end{proof}
        Recall that $\Hom_R(R, A) \cong A$ via the map $\cph \mapsto \cph(1)$ for any $R$-module $A$. Notice that 
        \begin{align*}
            d_1(\Hom_{\Z G}(\Z G, A)) &= \SET{\cph((\sigma-1)x) \mid \cph \in \Hom_{\Z G}(\Z G, A)} 
            \\&= \SET{(\sigma-1)\cph \mid \cph \in \Hom_{\Z G}(\Z G, A)} = (\sigma-1)\Hom_{\Z G}(\Z G, A)
        \end{align*}
        The image of the last term on the right under the above map is clearly $(\sigma-1)A$. Thus by the lemma we have that $H^1(G, A) \cong A/(\sigma-1)A$. Since the kernel of $d_n$ is 0 for $n \geq 3$, we see that $H^n(G, A) = 0$ for $n \geq 3$.
    \end{enumerate}

    \newpage
        \begin{theorem}
            $\Hom_R(\cdot, D)$ is a left-exact functor.
        \end{theorem}
        \begin{proof}
            I leave it as an exercise to see that $\Hom_R(\cdot, D)$ is a contravariant functor by sending the diagram
            \[\begin{tikzcd}
                A & B
                \arrow["\varphi", from=1-1, to=1-2]
            \end{tikzcd}\]
            to 
            \[\begin{tikzcd}
                {\Hom_R(B, D)} & {\Hom_R(A,D)}
                \arrow["{\psi \mapsto \psi \circ \varphi}", from=1-1, to=1-2]
            \end{tikzcd}\]

            We need to show that given a short exact sequence 
            \[\begin{tikzcd}
                A & B & C & 0
                \arrow["\psi", from=1-1, to=1-2]
                \arrow["\varphi", from=1-2, to=1-3]
                \arrow[from=1-3, to=1-4]
            \end{tikzcd}\]
            The sequence
            \[\begin{tikzcd}
                0 & {\Hom_R(C,D) } & {\Hom_R(B, D)} & {\Hom_R(A, D)}
                \arrow[from=1-1, to=1-2]
                \arrow["\tau", from=1-2, to=1-3]
                \arrow["\sigma", from=1-3, to=1-4]
            \end{tikzcd}\]
            is exact. If $\Psi \circ \varphi = \Psi' \circ \varphi$, since $\varphi$ is surjective it has a right inverse, which shows that $\Psi = \Psi'$. Given $\Phi \in \ker \sigma$, we know that $\Phi \circ \psi = 0$, so $\Im \psi \subset \ker \Phi$. Since $\Im \psi = \ker \cph$, we know that $\ker \cph \subset \ker \Phi$. We seek to find $a \in \Hom_R(C,D)$ so that $a \circ \cph = \Phi$. Thus, we look to define $a(c) = \Phi \circ \cph^{-1}(c)$. Notice that 
            \begin{align*}
                \cph(\cph^{-1}(rc)) = rc = \cph(r\cph^{-1}(c))
            \end{align*}
            This shows that $\cph^{-1}(rc) - r\cph^{-1}(c) \in \ker \cph \subset \Im \Phi$. Thus we see that $\Phi(\cph^{-1}(rc)) = \Phi(r\phi^{-1}(c)) = r\Phi(\phi^{-1}(c))$, so $\Phi \circ \cph^{-1}$ is indeed a $R$-module homomorphism, completing the proof.
        \end{proof}
\end{document}